\chapter{Classification, probability, and decisions}
\label{ch:classification}

\begin{learningobjectives}
After this chapter, you should be able to:
\begin{itemize}
  \item distinguish class scores, probability estimates, and decisions;
  \item derive logistic loss and its gradient;
  \item choose metrics and thresholds from a use-specific error structure; and
  \item evaluate discrimination, calibration, and subgroup behavior.
\end{itemize}
\end{learningobjectives}

\begin{chapterconnection}
Linear regression supplied the affine score $z=w^Tx+b$. Classification retains
that score but changes its interpretation, loss, and action rule. These choices
cannot be hidden inside a generic ``single neuron'' update without losing the
mathematical contract.
\end{chapterconnection}

\section{A score is not yet a probability or decision}

For binary $Y\in\{0,1\}$, an affine score orders examples along a direction.
The perceptron predicts from the sign and updates on classification mistakes.
It is historically and geometrically instructive, especially when data are
linearly separable, but it does not estimate a probability and has no finite
maximum-margin or likelihood solution merely because it converges.

Logistic regression maps the score through
\[
\sigma(z)=\frac{1}{1+e^{-z}},\qquad
\widehat p(x)=\sigma(w^Tx+b).
\]
The log-odds are affine:
\[
\log\frac{\widehat p(x)}{1-\widehat p(x)}=w^Tx+b.
\]
A probability estimate still does not specify an action. A policy chooses a
threshold, abstention region, or downstream optimization from costs and
constraints.

\section{Bernoulli likelihood and logistic loss}

Under a conditional Bernoulli model, the negative average log-likelihood is
\[
J(w,b)=-\frac{1}{n}\sum_{i=1}^n
\left[y_i\log p_i+(1-y_i)\log(1-p_i)\right].
\]
Its gradient with respect to $w$ is $X^T(p-y)/n$. The compact expression
resembles least squares, but it follows from a different loss-link pairing.
This is why a universal gradient formula attached to a generic neuron is
misleading.

Numerical implementations operate on logits with stable functions rather than
forming extreme probabilities and taking their logarithms. Regularization,
scaling, class weighting, and convergence status remain part of the fitted
object.

\conceptcheck{Why can rounding a predicted probability at 0.5 be both a valid
software operation and an invalid decision policy? Name the missing information.}

\section{The confusion matrix is conditional on a threshold}

At a chosen threshold, predictions produce true positives, false positives,
true negatives, and false negatives. Accuracy weights every row equally and can
look excellent when the event is rare. Precision asks what fraction of positive
predictions are positive; recall asks what fraction of positive outcomes are
detected. Specificity concerns negative outcomes.

No metric is universally best. The metric must represent a loss, utility, or
scientific purpose. Receiver-operating characteristic curves compare true- and
false-positive rates across thresholds. Precision-recall curves often expose
rare-event behavior more directly because precision depends on prevalence.
Area summaries discard the operating threshold and can hide unacceptable
regions.

\section{Thresholds belong to a decision policy}

If a false negative costs $C_{FN}$ and a false positive costs $C_{FP}$, an
idealized calibrated probability model with no other constraints acts when its
expected cost is lower. Real policies may also include capacity limits, review
queues, abstention, subgroup constraints, or asymmetric downstream actions.

Version the threshold separately from model parameters. It may change as costs,
prevalence, or operational capacity change without retraining the score model.
Evaluate the final model-policy pair on data untouched by threshold selection.

\begin{professionalbox}
An API should not label a floating-point score as ``risk'' without specifying
its semantics. Return a documented score or probability, model version, and
applicable population. If the service also returns an action, identify the
decision-policy version.
\end{professionalbox}

\section{Calibration asks whether probabilities mean what they say}

A model is calibrated on a population if events assigned probability near
$q$ occur at frequency near $q$. Discrimination and calibration differ. A
model can rank examples well while systematically overstating probability.
Reliability diagrams and proper scoring rules provide evidence, but bin choices
and sample size matter.

Calibration can shift when prevalence or measurement changes. Platt-style
scaling or isotonic calibration must be fitted on data separate from the data
used to fit the underlying score, commonly through nested or cross-validated
procedures.

\begin{misconceptionbox}
\textbf{Misconception:} the output of a sigmoid is automatically a trustworthy
probability. \textbf{Correction:} the range is between zero and one, but
probability meaning depends on the model, training objective, data-generating
setting, and empirical calibration on the target population.
\end{misconceptionbox}

\section{Multiclass problems require an explicit structure}

For mutually exclusive classes, a softmax transforms logits $z_k$ into
$p_k=e^{z_k}/\sum_j e^{z_j}$. Cross-entropy fits the probability assigned to the
observed class. One-vs-rest constructs separate binary problems and may require
additional reconciliation. Multilabel prediction is different: several labels
may be true simultaneously.

Report macro summaries when each class should receive equal weight and micro
summaries when each decision receives equal weight. Always inspect class-level
errors and the confusion matrix; a single average can conceal a class that the
system never recognizes.

\begin{workedexample}{degradation score to review decision}
A logistic model estimates whether a cell will cross a degradation threshold
within 100 cycles. The validation set is grouped by cell and later in time than
training. The laboratory can manually review only 30 cells per week.

The team first evaluates probability calibration. It then chooses a threshold
that respects review capacity while placing a stated penalty on missed events.
The held-out report contains prevalence, precision, recall, calibration, and
errors by manufacturing batch. The artifact stores the score model; a separate
policy file stores the threshold and capacity assumption. Neither output is
described as a safety certification.
\end{workedexample}

\section{Exercises}
\begin{enumerate}
  \item Derive the logistic-loss gradient from the Bernoulli likelihood.
  \item Construct two classifiers with the same accuracy but different false-
  negative costs for a stated application.
  \item Explain how a prevalence change can alter precision while leaving class-
  conditional score distributions unchanged.
  \item Specify tests for score range, class ordering, threshold boundaries,
  missing features, and policy-version compatibility.
\end{enumerate}

\begin{chapterreview}
Classification has at least three layers: a fitted score, a probability
interpretation, and a decision policy. Logistic regression supplies a principled
loss and conditional probability model; threshold metrics describe one action
rule; calibration audits probability meaning. Consequential behavior depends on
population, costs, capacity, and feedback, not on a default threshold.
\end{chapterreview}

\begin{notebooklink}
Lecture 11 notebook 00 connects perceptron geometry to logistic estimation.
Companion laboratories develop thresholds, calibration, multiclass structure,
and a versioned classification policy.
\end{notebooklink}
